\section{Introduction}
\label{iv:introduccion}
\begingroup
\fontsize{10.5}{12.8}\selectfont
\setlength{\parskip}{1pt plus .5pt minus .25pt}

Comment une même structure discrète peut-elle déterminer des coordonnées du
continu, soutenir une réalisation exceptionnelle et conserver une dualité
lors du passage à des représentations dimensionnelles ? La question exige plus
que de réunir des valeurs, des dimensions et des groupes : elle exige de construire les objets
intermédiaires et les applications qui transmettent leur information. Cet
article développe cette connexion dans l’holographie modulaire triadique (HMT).
Son résultat structurant est un double prolongement de l’état
généalogique : l’incidence conduit à une algèbre d’opérateurs de vertex
isomorphe à \(V^\natural\), tandis que les représentations de phase, d’action
et de canaux déterminent une dualité compacte et un morphisme d’écrans.
La coordination des deux prolongements précise ensuite le domaine
algébrique et les conditions de transport vers une réalisation de
théorie M.

La construction commence avant l’évaluation réelle.
\emph{All Possible Paths} (APP) désigne les chemins sur le support
\(X_9=(\mathbb Z/9\mathbb Z)^2\). Son graphe de Cayley utilise la loi
additive et les générateurs cardinaux \(\pm(1,0),\pm(0,1)\). Sur les
mêmes positions, on évalue la somme et le produit ; la réduction résiduelle
conserve le quotient de chaque évaluation. Les feuillets additif et
multiplicatif sont donc deux évaluations d’un support commun.
La signature \(\mathrm{TRIT}=\{-1,0,+1\}\) distingue l’orientation du
transport et les régimes hyperbolique, parabolique et elliptique de sa
réalisation quadratique. Le \emph{Topological Phase Kernel} (TPK) est la
composition effective de sélection d’évaluation, de transport cardinal
et de mise à jour des registres qui prolonge des états compatibles. Ces
états retiennent les positions, les feuillets, l’orientation, les résidus, les quotients,
les retenues, les préfixes, le bord et la mémoire. La sélection ne remplace pas le
déplacement, et le retour d’une phase n’efface pas l’histoire parcourue.

% BEGIN R27_FRAMING_INTRO_ES
Les opérations qui relient le support arithmétique à l’incidence
admettent des preuves intermédiaires explicites. Le changement de module, la répétition
du support à module fixe et la variation dimensionnelle ont leurs propres
lois du défaut somme--produit. La compression en trois classes produit
un sous-réseau d’indice deux et une unité quadratique de norme un.
Dans l’émetteur, l’alphabet décimal, les raccordements de demi-émissions et la
comparaison entre histogramme et mémoire ordonnée conservent les différences
entre valeur, incidence et trajectoire. Le choix de l’observable
multiplicatif est examiné avant la sélection régionale.

% BEGIN R28_FRAMING_INTRO_ES
Le multigraphe étiqueté de l’émetteur et son quotient de phase situent ces
réalisations sur des incidences explicites. Les cinq clôtures minimales
à huit arêtes couvrent les ancrages APP, de propagation, d’auto-échelle
et de chaque région de clôture. L’occupation minimale du bord et la
parité de ses lecteurs de mémoire précisent, à leur tour, la sélection des
représentants et la transformation des registres lors de l’inversion du chemin.
% END R28_FRAMING_INTRO_ES

Le premier résultat est l’organisation conjointe de la structure
discrète du continu. Dans la section~\ref{iv:continuo-conjunto}, un
même arbre d’exécutions détermine la mesure cylindrique et le transport
des feuillets, le bilan orienté, l’incidence ternaire, le complexe
de mémoire et son assemblage terminal avec droite déterminante. Les applications
sont définies sur les générateurs et leur compatibilité avec les
restrictions d’histoire et les réductions modulaires est prouvée. La mesure
utilise des poids positifs compatibles dans un arbre à ramification finie
sans feuilles terminales ; l’identité de Gram conserve le résidu terminal
de survie ; la construction déterminante maintient ses degrés et
ses bases, avec les hypothèses supplémentaires qu’exige une évaluation
scalaire. Les cinq constructions conservent ainsi une même origine et les
conditions particulières de chaque opération. Ce ne sont pas cinq domaines
sélectionnés après la production des coordonnées.

Les sections de cylindres compatibles déterminent les coordonnées
numériques de la génération corrélée de \(\pi,\varphi,e,\alpha\)
sous les conditions d’emboîtement, de contraction et de bord
de leur réalisation archimédienne. La projection exceptionnelle conserve,
en revanche, les données d’incidence réalisées comme des codes,
des réseaux et des algèbres graduées. Les feuillets surfacique et volumétrique et leurs
mesures, développés dans la section~\ref{sec:medida-retorno-hojas},
et les vacances associées au raffinement, dans la
section~\ref{vac-capacidad-trit}, interviennent avec les registres
régionaux dans la construction de \(\alpha\) et dans le prolongement
exceptionnel. La masse électronique utilise une autre composition de ces
opérations communes ; son évaluation n’est pas l’objet de cet article.

Le système de coordonnées qui regroupe deux trits en un chiffre nonadique dispose d’un inverse
et d’une somme transportée avec retenue ; la projection du cube ambiant
à \(729\) mots sur ses deux premières coordonnées a des fibres
uniformes de neuf éléments. La preuve de corrélation quadratique
détermine ensuite l’identité du système de coordonnées à six positions.
Sur le bord, la saturation laisse huit possibilités, la neutralité
duale en laisse deux et l’orientation sélectionne le représentant fixé.
Les contrôles de clôture conservent les extrémités de l’identité
télescopique. 

La connexion exceptionnelle a un antécédent concret dans le registre
dodécaphasique \(K\). Ses douze composantes proviennent du registre signé
et des opérations du TPK. Le codage rationnel conserve les
blocs lorsque la base, la longueur, l’origine de phase et l’orientation sont fixées ;
l’inversion de Hadamard récupère la coordonnée signée correspondante.
Les supports du registre interviennent dans un drapeau marqué du plan
de Witt. À partir du code ternaire, on construit le recollement du réseau
de Niemeier de racines \(A_2^{12}\) ; le drapeau et l’origine, avec la
métrique et la chambre positive déclarées, sélectionnent le voisin sans
racines qui est identifié au réseau de Leech. Chaque opération
transmet une donnée nécessaire à l’objet suivant.

% BEGIN S27_FRAMING
La section~\ref{star27:comparacion} précise le passage entre les registres
régionaux et l’incidence par la comparaison de trois étoiles.
Les intersections de \(P_\pi\) avec \(H_e\), \(H_\varphi\) et
\(H_{\rm APP}\) déterminent les quadruplets \(B_K\), \(B_\varphi\) et
\(B_{\rm APP}\) ; la partition de signe de la cochaîne antérieure à la
retenue détermine la polarisation de leurs complétions. Le recensement des
\(495\) quadruplets situe ces configurations dans le plan complet,
et le stabilisateur du repère ordonné de quatre supports caractérise
sa rigidité dans \(G_W\). Cette rigidité concerne l’action sur les
positions ; la détermination numérique de \(K\) conserve sa construction
entière, l’origine de phase et l’orientation préalablement fixées.

La section~\ref{carry27:seccion} développe la composante entière qui
accompagne la lecture résiduelle. La division euclidienne produit les
résidus \(q,a\) et les quotients \(c,h\), dont la composition conserve la
mémoire des multiples de trois. Le cocycle, les représentations
polynomiales sur \(\mathbb F_3\) et les identités de transport
explicitent quelles données traversent chaque changement de registre. Leur lien avec
les cylindres compatibles conserve les bornes de bord et les fibres
de la lecture régionale. La comparaison des étoiles et la loi de
retenue remplissent ainsi des fonctions d’incidence et numériques déterminées
dans la connexion que cet article prolonge vers les réseaux et les algèbres.

% END S27_FRAMING


Sur ce réseau, on construit l’algèbre réticulaire, avec oscillateurs,
extension centrale et produit de vertex. L’extension d’ordre deux de
Frenkel--Lepowsky--Meurman et l’extension d’ordre trois associée à
l’isométrie ternaire réalisent \(V^\natural\), une fois vérifiées les
hypothèses des théorèmes de reconnaissance utilisés
\cite{flm,chenlamshimakura2016,abelamyamada2017}.
Les isomorphismes choisis entre les deux présentations conservent le vide,
le vecteur conforme, la graduation et les produits, et transportent les automorphismes
et les traces. Le Monstre est le groupe des automorphismes de cette algèbre ;
son action a pour domaine \(V^\natural\), non les chiffres de
\(\alpha\). Le théorème de Moonshine de Borcherds détermine les traces
graduées correspondantes~\cite{borcherds}. La construction du
réseau, l’extension algébrique et la reconnaissance de ces traces
sont des étapes distinctes d’une même chaîne d’applications.

Lorsque les traces de niveau trois sont atteintes, la factorisation formelle
de l’identité rationnelle détermine une divisibilité valable à tous
les degrés, avec son exposant optimal.
% END R27_FRAMING_INTRO_ES

Le registre \(K\) possède en outre une réalisation directionnelle.
Dans la représentation déclarée de \(A_5\) sur les douze positions,
sa projection sur le secteur tridimensionnel détermine un vecteur
\(u_K\ne0\) dans \(V_{11}=\mathbf1^\perp\).
Le vecteur non nul détermine la droite \(\ell_K\), et sa norme quadratique
exacte normalise le terme de rang un du projecteur sur
\(V_{10}=V_{11}\cap u_K^\perp\), dont le noyau et la covariance sont
spécifiés. La réduction \(12\to11\) élimine le mode uniforme ;
la réduction \(11\to10\) élimine la direction sélectionnée par le
registre complet. D’autre part, son orbite cyclique est une référence
inversible pour identifier des opérateurs au moyen de leurs réponses.
La section~\ref{sec:acoplamiento-recuperacion} établit les bornes
de stabilité et le bilan bilatéral : deux canaux construits à partir d’isométries permettent
de récupérer l’état ; l’extension unitaire inclut la mémoire entrante,
et la conservation d’un observable requiert la condition de commutation
de son énoncé.

La réalisation hilbertienne fournit le support opératoriel de la
dualité. À chaque profondeur, les phases étiquettent une base orthonormale
et les opérateurs de décalage et de caractère engendrent toute son algèbre
matricielle. Leur relation de Weyl représente l’extension centrale des
translations avec le signe fixé de la forme d’aire orientée.
L’inclusion au moyen de l’identité de toute la nouvelle fibre conserve
l’unité et la trace normalisée : la limite est l’algèbre uniformément
hyperfinie nonadique et sa clôture traciale est un facteur de type \(II_1\).
Fixer une continuation produit, en revanche, une inclusion dans un coin
et une autre clôture. Le choix de l’application de raffinement détermine donc
le type de réalisation et pas seulement sa dimension.

La section d’action et les coordonnées angulaires fixent le rayon qui
spécialise la dualité compacte. La norme déterminantale, les
coefficients du caractère d’action et le renouvellement régional normalisé
déterminent les sections positives antérieure et postérieure au retour.
Les systèmes de coordonnées angulaires conservent l’orientation et le nombre de tours.
Leur composition définit une ellipse étiquetée dont le produit des demi-axes
conserve l’action et dont le quotient détermine un rayon sans dimension
unique. Les conditions de positivité situent les coordonnées engendrées
dans le système de coordonnées elliptique ; l’échange des demi-axes produit le rayon
réciproque en conservant l’échelle commune d’action.

Sur \(\mathbb Z^2\times\mathbb R_{>0}\), l’échange des
coordonnées entières et l’inversion du rayon forment une involution
qui conserve la forme quadratique compacte. Le Hamiltonien diagonal est
autoadjoint, non négatif et à résolvante compacte ; la permutation des
coordonnées le conjugue unitairement avec le Hamiltonien de rayon
réciproque, y compris les domaines des opérateurs et des formes quadratiques.
Le changement de section résiduelle transporte aussi la relation \(L_9\)
qui conserve l’entier. L’énergie minimale de chaque classe est bien
définie et est évaluée au moyen de deux candidats entiers. La normalisation
par l’échelle de corde exprime la même dualité en rayons
dimensionnels. Sur l’axe imaginaire du demi-plan supérieur, l’inversion
modulaire lui est reliée par une semi-conjugaison explicite
de paramètres. Cette application n’identifie pas l’opérateur d’échange aux
automorphismes des secteurs d’un orbifold.

Les écrans dimensionnels coordonnent la direction \(K\), l’incidence
et le secteur compact. Le poinçonnage du code ternaire produit un
code de longueur onze avec distance cinq et partition parfaite.
La réduction orthogonale \(12\to11\to10\) et la réduction réticulaire
isotrope \(26\to24\) conservent leurs propres présentations, relations
et noyaux. Le morphisme \(\mathcal R_M^{\mathrm{alg}}\) identifie leurs
présentations quotientées et entrelace la dualité ; le graphe de la
projection conserve la coordonnée de \(V_{11}\) avec son image
dans \(V_{10}\). Les mêmes opérations se restreignent à la collection des
deux rayons déterminés par l’ellipse.

Le prolongement dynamique commence par la composition des trajectoires,
non par le choix des champs. En conservant la donnée précédente au
bord, les changements de direction et de type arithmétique définissent un
cocycle additif d’action. La surface d’univers et l’action de
Polyakov spécifient ensuite une réalisation variationnelle. L’application
vers la réalisation de théorie M reçoit les écrans dans un codomaine
à onze dimensions avec métrique, forme de degré trois, gravitino, graduation,
action et contraintes. Sous les hypothèses d’entrelacement et de
transport des domaines, le théorème de supercharge conserve \(Q^2=H\)
sur l’image de l’isométrie. L’identification physique complète
requiert en outre les tensions, les oscillateurs et les amplitudes, le contrôle
des anomalies et les limites déclarées dans la
section~\ref{iv:condiciones-realizacion-fisica}.

Cette hiérarchie détermine le développement : noyau et continu conjoint ;
génération coinductive, registre \(K\) et les deux constructions de
\(\alpha\) ; incidence et algèbre exceptionnelle ; réalisation hilbertienne,
action et rayon ; dualité, écrans et transport dynamique.
La section~\ref{sec:k-registro} détermine le registre utilisé dans
\ref{exc:k-direccion} ; la réalisation de \ref{iv:hilbert} soutient
les opérateurs de \ref{iv:dualidad} ; et
\ref{iv:pantallas} compose les réductions dans le morphisme que reçoit la
réalisation à onze dimensions. L’argument relie ainsi la structure,
la valeur, la représentation et l’équation au moyen de leurs applications effectives.

\par\endgroup
